\documentclass[11pt,a4paper]{article}
\usepackage[T1]{fontenc}
\usepackage{lmodern}
\usepackage{microtype}
\usepackage[margin=27mm]{geometry}
\usepackage{amsmath,amssymb,amsthm,mathtools}
\usepackage{xcolor}
\usepackage{enumitem}
\usepackage[hypertexnames=false]{hyperref}
\hypersetup{colorlinks=true,linkcolor=blue!45!black,urlcolor=blue!45!black,
  pdftitle={A single NAND axiom for Boolean algebra},
  pdfsubject={An explicit equational proof and reconstruction of the Boolean lattice}}
\newtheorem{theorem}{Theorem}
\newtheorem{lemma}[theorem]{Lemma}
\newcommand{\one}{\mathbf{1}}
\newcommand{\zero}{\mathbf{0}}
\newcommand{\W}{\mathrm{W}}
\setlength{\emergencystretch}{2em}
\setlist[enumerate]{itemsep=2pt,topsep=5pt}
\allowdisplaybreaks[1]
\title{A single NAND axiom for Boolean algebra}
\author{}
\date{}

\begin{document}
\maketitle
\vspace{-2em}
\begin{abstract}
We give an explicit proof that the equation
\[
 ((a*b)*c)*\bigl(a*((a*c)*a)\bigr)=c
\]
characterizes the NAND reducts of Boolean algebras on nonempty carriers.
The proof has two parts: a sequence of substitutions establishes
commutativity and double negation; an order-theoretic construction then
recovers the Boolean lattice. We also explain why the equation is complete
for identities in the NAND language.
\end{abstract}

\section{Statement and plan}
Let $A$ be a nonempty set equipped with a binary operation $*$ satisfying
\begin{equation}
 \boxed{((x*y)*z)*\bigl(x*((x*z)*x)\bigr)=z}
 \tag{W}\label{eq:W}
\end{equation}
for all $x,y,z\in A$. No other property of $*$ is assumed. In particular,
parentheses must be retained: even in a Boolean algebra, NAND is not
associative.

\begin{theorem}\label{thm:main}
There is a Boolean algebra structure on $A$ whose NAND operation is $*$.
Writing $n(x)=x*x$, its operations are
\[
 \neg x=n(x),\qquad
 x\wedge y=n(x*y),\qquad
 x\vee y=n(x)*n(y),
\]
and its constants are
\[
 \one=x*n(x),\qquad \zero=n(\one).
\]
The expression defining $\one$ is independent of $x$. Conversely, NAND
in any Boolean algebra satisfies \eqref{eq:W}.
\end{theorem}

We include the one-element Boolean algebra. Nonemptiness is needed to
obtain the constants: an empty algebra satisfies \eqref{eq:W} vacuously.

The difficult part is the following purely equational lemma. Its proof
occupies Section~\ref{sec:equations}. Section~\ref{sec:lattice} explains
how its identities determine a bounded distributive lattice with
complement, and Section~\ref{sec:complete} proves the converse and
completeness.

\begin{lemma}[The algebraic lemma]\label{lem:core}
Put $n(x)=x*x$. The single identity \eqref{eq:W} implies
\begin{align}
 x*y&=y*x,\tag{C}\label{eq:comm}\\
 n(n(x))&=x,\tag{DN}\label{eq:DN}\\
 n(x)*(x*y)&=x,\tag{I}\label{eq:I}\\
 (x*y)*(y*(x*z))&=y,\tag{S}\label{eq:S}\\
 x*(y*(x*(y*z)))&=x*(y*z).\tag{J}\label{eq:J}
\end{align}
\end{lemma}

\section{From the single axiom to the algebraic lemma}\label{sec:equations}
The repeated right-hand factor in \eqref{eq:W} suggests introducing
\[
 q(x,y)=x*((x*y)*x),\qquad m(x)=q(x,x).
\]
Then the axiom reads
\[
 ((x*y)*z)*q(x,z)=z.
\]
We write $\W(x,y,z)$ for this instance of the axiom. Our intermediate
goals are $m(x)=x*x$ and $q(x,y)=y*x$. These simplify the repeated
expressions enough to reveal double negation and commutativity.
Until commutativity has been proved, the order of every pair of factors
is significant.

\subsection{Four substitution identities}
The instance $\W(x,(x*y)*x,y)$ gives
\begin{equation}
 (q(x,y)*y)*q(x,y)=y.\tag{R}\label{eq:R}
\end{equation}
Expanding the outer $q$ and applying \eqref{eq:R}, we obtain
\begin{equation}
 q(q(x,y),y)=q(x,y)*y.\tag{R'}\label{eq:Rprime}
\end{equation}

For the next identity, use
$\W((v*w)*x,q(v,x),y)$. Its left-hand side is
\[
 \Bigl(\bigl(((v*w)*x)*q(v,x)\bigr)*y\Bigr)*q((v*w)*x,y).
\]
The parenthesized factor equals $x$ by $\W(v,w,x)$, so
\[
 (x*y)*q((v*w)*x,y)=y.
\]
Every element $z$ can be written as a product, since $\W(z,z,z)$ gives
$z=((z*z)*z)*m(z)$. Thus $v*w$ in the preceding identity can be any
$z\in A$, and
\begin{equation}
 (x*y)*q(z*x,y)=y.\tag{$C_0$}\label{eq:C0}
\end{equation}

In \eqref{eq:C0}, substitute $x=q(a,b)$, $y=c$, and $z=(a*a)*b$.
The instance $\W(a,a,b)$ simplifies $z*x$ to $b$. Hence
\begin{equation}
 (q(a,b)*c)*q(b,c)=c.\tag{$D_0$}\label{eq:D0}
\end{equation}
Finally, substitute $(a,b,c)=(x,q(y,z),z)$ into \eqref{eq:D0} and
use \eqref{eq:Rprime}:
\begin{equation}
 (q(x,q(y,z))*z)*(q(y,z)*z)=z.\tag{$J_0$}\label{eq:J0}
\end{equation}

\subsection{Identifying the square and proving double negation}
Fix $a\in A$ and abbreviate
\[
 u=m(a),\qquad t=(a*a)*a.
\]
The definition of $m$ and $\W(a,a,a)$ give
\begin{equation}
 a*t=u,\qquad t*u=a.\label{eq:atu}
\end{equation}
We first derive four relations involving $a$ and $u$:
\begin{align}
 (a*u)*a&=u,\tag{$U_1$}\label{eq:U1}\\
 u*u&=a,\tag{$U_2$}\label{eq:U2}\\
 u*(u*a)&=a,\tag{$U_3$}\label{eq:U3}\\
 a*(u*(a*u))&=u.\tag{$U_4$}\label{eq:U4}
\end{align}
To prove \eqref{eq:U1}, apply \eqref{eq:C0} with $(x,y,z)=(a,u,a*a)$.
This gives $(a*u)*q(t,u)=u$, while
\[
 q(t,u)=t*((t*u)*t)=t*(a*t)=t*u=a.
\]
Now $\W(a,u,a)$ becomes $((a*u)*a)*u=a$; applying
\eqref{eq:U1} proves \eqref{eq:U2}.

For \eqref{eq:U3}, note that
$q(a,u)=a*((a*u)*a)=a*u$. In \eqref{eq:J0}, set $x=y=z=a$.
Since $q(a,a)=u$, the result is
\[
 ((a*u)*a)*(u*a)=a,
\]
which reduces to \eqref{eq:U3}. Finally, \eqref{eq:C0} with
$(x,y,z)=(t,u,a)$ gives $a*q(u,u)=u$. Expanding $q(u,u)$ and using
$u*u=a$ proves \eqref{eq:U4}.

We need one identity that applies these relations at a different element.
For arbitrary $d,b$, set
\[
 v=(d*d)*d,\qquad r=q(b*(d*d),d).
\]
By \eqref{eq:C0}, $v*r=d$. Therefore
\[
 q(v,r)=v*((v*r)*v)=v*(d*v)=v*m(d)=d,
\]
where the last equality is $\W(d,d,d)$. Now \eqref{eq:J0} at
$(x,y,z)=(v,b*(d*d),d)$ yields
\[
 (d*d)*(q(b*(d*d),d)*d)=d.
\]
Replace $d$ by $m(a)$ and use \eqref{eq:U2}, which holds for every $a$:
\begin{equation}
 a*(q(b*a,m(a))*m(a))=m(a).\tag{V}\label{eq:V}
\end{equation}

Return to our fixed $a$ and put $w=u*(a*u)$. By \eqref{eq:U4}
and \eqref{eq:U2}, respectively,
\begin{equation}
 a*w=u,\qquad m(u)=u*((u*u)*u)=w.\label{eq:w}
\end{equation}
The instance $\W(u,u,u*a)$ reduces, using \eqref{eq:U2} and
\eqref{eq:U3}, to
\begin{equation}
 (a*(u*a))*w=u*a.\label{eq:waux}
\end{equation}
Apply \eqref{eq:V} with $(a,b)$ replaced by $(u,t)$. Since $t*u=a$
and $m(u)=w$, this gives
\[
 u*(q(a,w)*w)=w.
\]
But $q(a,w)=a*((a*w)*a)=a*(u*a)$. By \eqref{eq:waux}, the last
displayed left-hand side is $u*(u*a)=a$. Thus $w=a$, and $a*w=u$
now gives
\begin{equation}
 \boxed{m(a)=a*a.}\label{eq:m}
\end{equation}
Substituting this into $u*u=a$ proves double negation \eqref{eq:DN}.
We also record two consequences needed below:
\begin{align}
 (a*a)*(a*(a*a))&=a,\tag{H}\label{eq:H}\\
 ((a*b)*a)*(a*a)&=a.\tag{L}\label{eq:L}
\end{align}
The first is $w=a$ after substituting $u=a*a$; the second is
$\W(a,b,a)$ after substituting $m(a)=a*a$.

\subsection{Identifying the reversed product}
We next prepare identities that let us compare $q(a,b)$ with $b*a$.
First, $\W(q(a,b),c,b)$ and \eqref{eq:Rprime} give
\begin{equation}
 ((q(a,b)*c)*b)*(q(a,b)*b)=b.\tag{$P_0$}\label{eq:P0}
\end{equation}
Apply \eqref{eq:C0} with
\[
 x=q(a,b)*b,\qquad y=c,\qquad z=(q(a,b)*b)*b.
\]
By \eqref{eq:P0} at $c=b$, the product $z*x$ is $b$, so
\begin{equation}
 ((q(a,b)*b)*c)*q(b,c)=c.\tag{$E_0$}\label{eq:E0}
\end{equation}
In this identity replace $(a,b,c)$ by $(b,a,q(b,a))$.
Its first factor is $a$ by \eqref{eq:R}, giving
\begin{equation}
 a*q(a,q(b,a))=q(b,a).\tag{$K_0$}\label{eq:K0}
\end{equation}
Two more substitutions into the original axiom yield
\begin{align}
 a*q(b*c,q(b,a))&=q(b,a),\tag{$B_0$}\label{eq:B0}\\
 a*q(b,q(c*b,a))&=q(c*b,a).\tag{$I_0$}\label{eq:I0}
\end{align}
For \eqref{eq:B0}, use $\W(b*c,a,q(b,a))$ and simplify the first
factor by $\W(b,c,a)$. For \eqref{eq:I0}, use
$\W(b,a,q(c*b,a))$ and simplify the first factor by \eqref{eq:C0}.

We first identify $q$ on pairs whose second entry is a square:
\begin{equation}
 \boxed{q(b*a,a*a)=a.}\tag{$T_0$}\label{eq:T0}
\end{equation}
Here are the intermediate steps. Let $r=q(x*y,y)$. Equation \eqref{eq:C0}
gives $(y*y)*r=y$, and hence
\[
 q(y*y,r)=(y*y)*(y*(y*y))=y
\]
by \eqref{eq:H}. In \eqref{eq:I0}, take $(a,b,c)=(r,y,y)$ to get
\[
 r*q(y,q(y*y,r))=q(y*y,r).
\]
This reduces to $r*q(y,y)=y$. Since $q(y,y)=m(y)=y*y$,
\begin{equation}
 q(x*y,y)*(y*y)=y.\tag{$N_0$}\label{eq:N0}
\end{equation}
Now apply the original axiom at
\[
 \W\bigl(b*a,\,((b*a)*a)*(b*a),\,a*a\bigr).
\]
The first factor is $q(b*a,a)*(a*a)=a$ by \eqref{eq:N0}, leaving
\begin{equation}
 a*q(b*a,a*a)=a*a.\tag{$O_0$}\label{eq:O0}
\end{equation}
Put $s=q(b*a,a*a)$. Equation \eqref{eq:I0} at $(a*a,a,b)$ says
$(a*a)*q(a,s)=s$. By \eqref{eq:O0}, $a*s=a*a$, so
\[
 q(a,s)=a*((a*s)*a)=a*((a*a)*a)=m(a)=a*a.
\]
Thus $s=(a*a)*(a*a)=a$, proving \eqref{eq:T0}.

\subsection{Commutativity}
Fix $a,b$ and write
\[
 t=(a*b)*a,\qquad r=q(a,b)=a*t.
\]
The next three equalities relate $a,t,r$:
\begin{equation}
 t*r=a,\qquad r*a=t,\qquad t*a=r.\label{eq:balance}
\end{equation}
Indeed, \eqref{eq:T0} and \eqref{eq:L} give
$q(t,a*a)=a$ and $t*(a*a)=a$. Expanding the former,
\[
 a=q(t,a*a)=t*((t*(a*a))*t)=t*(a*t)=t*r.
\]
This proves the universal identity
\[
 ((x*y)*x)*q(x,y)=x.
\]
Apply it with $(x,y)=(t,a*a)$. The two known simplifications give
$(a*t)*a=t$, or $r*a=t$.
Finally,
\[
 q(a,t)=a*((a*t)*a)=a*(r*a)=a*t=r,
\]
and
\[
 q(t,r)=t*((t*r)*t)=t*(a*t)=t*r=a.
\]
Equation \eqref{eq:K0} with $(a,b)=(t,a)$ is therefore $t*a=r$.
This proves all of \eqref{eq:balance}.

We now establish
\begin{equation}
 (a*a)*q(a,b)=a.\tag{$W_1$}\label{eq:W1}
\end{equation}
For an arbitrary $c$, \eqref{eq:C0} gives $(a*a)*q(c*a,a)=a$.
Consequently,
\[
 q(a*a,q(c*a,a))=(a*a)*(a*(a*a))=a
\]
by \eqref{eq:H}. Thus \eqref{eq:J0} at $(x,y,z)=(a*a,c*a,a)$ becomes
\begin{equation}
 (a*a)*(q(c*a,a)*a)=a.\label{eq:W1aux}
\end{equation}
Take $c=t$. Using \eqref{eq:balance},
\[
 q(t*a,a)=q(r,a)=r*((r*a)*r)=r*(t*r)=r*a=t.
\]
Equation \eqref{eq:W1aux} becomes $(a*a)*(t*a)=a$; since $t*a=r$,
this is \eqref{eq:W1}.

We can now identify $q$ everywhere. Equation \eqref{eq:B0} at $c=b$ says
\[
 a*q(b*b,q(b,a))=q(b,a).
\]
Its inner expression simplifies by \eqref{eq:W1} and \eqref{eq:H}:
\[
 \begin{aligned}
 q(b*b,q(b,a))
   &=(b*b)*(((b*b)*q(b,a))*(b*b))\\
   &=(b*b)*(b*(b*b))=b.
 \end{aligned}
\]
We have proved
\begin{equation}
 \boxed{q(b,a)=a*b.}\tag{$Q_1$}\label{eq:Q1}
\end{equation}
Finally, rewrite \eqref{eq:K0} using \eqref{eq:Q1}:
\[
 \begin{aligned}
 q(b,a)&=a*q(a,q(b,a))\\
       &=a*(q(b,a)*a)\\
       &=a*((a*b)*a)\\
       &=q(a,b).
 \end{aligned}
\]
The first and last expressions are $a*b$ and $b*a$, respectively.
This proves commutativity \eqref{eq:comm}.

\subsection{The remaining three identities}
We may now use commutativity freely. Equation \eqref{eq:W1}, together
with \eqref{eq:Q1}, gives $n(x)*(x*y)=x$, proving \eqref{eq:I}.
The original axiom and \eqref{eq:Q1} give
\[
 ((x*y)*z)*(x*z)=z.
\]
Replacing $(y,z)$ by $(z,y)$ and commuting the indicated factors yields
$(x*y)*(y*(x*z))=y$, which is \eqref{eq:S}.
Finally, \eqref{eq:I0} and \eqref{eq:Q1} give
\[
 a*((a*(c*b))*b)=a*(c*b).
\]
Commutativity turns this into \eqref{eq:J}.
Together with double negation, this finishes the proof of
Lemma~\ref{lem:core}.\hfill$\square$

\section{Recovering the Boolean lattice}\label{sec:lattice}
From this point onward we use only the five identities in
Lemma~\ref{lem:core}. In particular, $*$ is commutative and $n$ is an
involution, so $n$ is both injective and surjective.

\subsection{Absorption identities and constants}
We begin with three useful consequences:
\begin{align}
 x*(x*y)&=x*n(y),\tag{B}\label{eq:B}\\
 x*(n(x)*y)&=n(x),\tag{I'}\label{eq:Iprime}\\
 x*n(x*y)&=x*y.\tag{A}\label{eq:A}
\end{align}
For \eqref{eq:B}, put $p=x*y$. Identity \eqref{eq:J} at $(x,p,y)$ is
\[
 x*(p*(x*(p*y)))=x*(p*y).
\]
Identity \eqref{eq:S} at $(y,x,p)$ gives $p*(x*(p*y))=x$.
Consequently $x*(y*(x*y))=n(x)$. Use this last equality with $x,y$
interchanged to simplify \eqref{eq:J} at $(x,y,x)$:
\[
 x*n(y)=x*(x*y).
\]
This is \eqref{eq:B}.

Replacing $x$ by $n(x)$ in \eqref{eq:I} proves \eqref{eq:Iprime}.
For \eqref{eq:A}, apply \eqref{eq:I} at $(x*y,n(x))$. Its inner
product is $(x*y)*n(x)=x$ by \eqref{eq:I} and commutativity, so
$n(x*y)*x=x*y$, as required.

Let $e(x)=x*n(x)$. For $d=x*y$, the preceding identities give
\[
 \begin{aligned}
 e(x)&=n(x)*x\\
     &=n(x)*(n(x)*d) &&\text{by \eqref{eq:I}}\\
     &=n(x)*n(d) &&\text{by \eqref{eq:B}}\\
     &=n(d)*(n(d)*x) &&\text{by \eqref{eq:B} and commutativity}\\
     &=n(d)*d &&\text{by \eqref{eq:A} and commutativity}\\
     &=e(d).
 \end{aligned}
\]
Interchanging $x,y$ gives $e(y)=e(d)$ as well. Thus $e(x)$ is constant,
and nonemptiness allows us to define
\[
 \one=e(x),\qquad \zero=n(\one).
\]
Equation \eqref{eq:B} and double negation show that
\begin{equation}
 x*\one=x*(x*n(x))=x*n(n(x))=n(x).\label{eq:one}
\end{equation}
Applying \eqref{eq:I} at $(\one,x)$ gives
$\zero*(\one*x)=\one$, hence $\zero*n(x)=\one$.
Since $n$ is surjective,
\begin{equation}
 \zero*x=\one\qquad\text{for all }x.\label{eq:zero}
\end{equation}

\subsection{The order and its reversal by negation}
Define a relation on $A$ by
\begin{equation}
 x\leq y\quad\Longleftrightarrow\quad x*y=n(x).\label{eq:order}
\end{equation}
It is reflexive by the definition of $n$. To see antisymmetry, suppose
$x\leq y$ and $y\leq x$. Commutativity gives $n(x)=x*y=y*x=n(y)$,
and injectivity of $n$ gives $x=y$.

A second description of the relation will be useful:
\begin{equation}
 x\leq y\quad\Longleftrightarrow\quad x*n(y)=\one.\label{eq:orderone}
\end{equation}
If $x*y=n(x)$, then \eqref{eq:B} gives
$x*n(y)=x*(x*y)=x*n(x)=\one$.
Conversely, if $x*n(y)=\one$, apply \eqref{eq:B} at $(x,n(y))$:
\[
 x*y=x*(x*n(y))=x*\one=n(x).
\]

For transitivity, suppose $x\leq y$ and $y\leq z$. Then
$y*x=n(x)$ and $y*n(z)=\one$. Identity \eqref{eq:S} at $(y,n(z),x)$,
\[
 (y*n(z))*(n(z)*(y*x))=n(z),
\]
therefore becomes
$\one*(n(z)*n(x))=n(z)$.
By \eqref{eq:one} this says $n(n(z)*n(x))=n(z)$.
Apply $n$ to both sides to obtain $n(x)*n(z)=z$.
Now \eqref{eq:Iprime} at $(x,n(z))$ gives
\[
 x*z=x*(n(x)*n(z))=n(x),
\]
so $x\leq z$. This proves that \eqref{eq:order} is a partial order.

Negation reverses it. Indeed, if $x\leq y$, then
\[
 n(y)*n(x)=n(y)*(y*x)=y=n(n(y))
\]
by \eqref{eq:I}, so
\begin{equation}
 x\leq y\quad\Longrightarrow\quad n(y)\leq n(x).\label{eq:anti}
\end{equation}
Since $n$ is involutive, this implication is an equivalence.

\subsection{Meet and join}
Define $x\wedge y=n(x*y)$. By \eqref{eq:A},
\[
 n(x*y)*x=x*y=n(n(x*y)).
\]
Thus $x\wedge y\leq x$, and commutativity also gives
$x\wedge y\leq y$.

We must prove that this lower bound is the greatest one. Suppose
$z\leq a$ and $z\leq b$, and put
\[
 k=n(z),\qquad p=a*b,\qquad h=k*p.
\]
The hypotheses say $a*z=b*z=k$. Applying \eqref{eq:S} at
$(a,t,z)$ and $(b,t,z)$ gives, for every $t$,
\begin{equation}
 (t*a)*(t*k)=t,\qquad (t*b)*(t*k)=t.\label{eq:loweraux}
\end{equation}
The second equality at $t=a$ is $p*(a*k)=a$.
Identity \eqref{eq:S} at $(k,p,a)$ now gives
\[
 h*a=(k*p)*(p*(k*a))=p.
\]
Use the first equality of \eqref{eq:loweraux} at $t=h$ to calculate
\[
 \begin{aligned}
 h&=(h*a)*(h*k)\\
  &=p*(k*(k*p))\\
  &=p*(k*n(p)) &&\text{by \eqref{eq:B}}\\
  &=n(p) &&\text{by \eqref{eq:Iprime} and commutativity}.
 \end{aligned}
\]
Finally, \eqref{eq:I} at $(k,p)$ is $n(k)*(k*p)=k$. Since
$n(k)=z$ and $k*p=h=n(p)$, it says
\[
 z*n(p)=n(z).
\]
This is exactly $z\leq n(p)=a\wedge b$. Hence $\wedge$ is the
greatest-lower-bound operation. Its commutativity, associativity,
idempotence, and monotonicity follow from this universal property.

An order-reversing involution turns greatest lower bounds into least
upper bounds. Explicitly, define
\[
 x\vee y=n\bigl(n(x)\wedge n(y)\bigr)=n(x)*n(y).
\]
Since $n(x)\wedge n(y)\leq n(x),n(y)$, applying $n$ shows
$x,y\leq x\vee y$. If $x,y\leq z$, then
$n(z)\leq n(x),n(y)$, so $n(z)\leq n(x)\wedge n(y)$;
applying $n$ gives $x\vee y\leq z$.
Thus $\vee$ is the least-upper-bound operation, and we have a lattice.
The same construction gives the De Morgan laws
\begin{equation}
 n(x\wedge y)=n(x)\vee n(y),\qquad
 n(x\vee y)=n(x)\wedge n(y).\label{eq:demorgan}
\end{equation}

By \eqref{eq:one}, $x\leq\one$. By \eqref{eq:zero} and double negation,
$\zero*x=\one=n(\zero)$, so $\zero\leq x$.
Finally,
\begin{equation}
 x\wedge n(x)=n(x*n(x))=\zero,\qquad
 x\vee n(x)=n(x)*x=\one.\label{eq:complement}
\end{equation}
We have obtained a bounded lattice in which $n(x)$ is a complement of $x$.

\subsection{Distributivity from splitting}
A complemented lattice need not be distributive. The additional information
we need comes from \eqref{eq:S} at $(x,y,z)=(a,c,a)$:
\[
 (a*c)*(c*n(a))=c.
\]
In our lattice notation this is the splitting identity
\begin{equation}
 \boxed{c=(c\wedge a)\vee(c\wedge n(a)).}\label{eq:split}
\end{equation}
Indeed, expanding the right-hand side as NAND terms and applying double
negation gives $(c*a)*(c*n(a))$.

Splitting first yields the useful criterion
\begin{equation}
 u\leq v\quad\Longleftrightarrow\quad u\wedge n(v)=\zero.
 \label{eq:orthogonal}
\end{equation}
For the forward direction, monotonicity of meet gives
$u\wedge n(v)\leq v\wedge n(v)=\zero$.
For the reverse direction, \eqref{eq:split} gives
\[
 u=(u\wedge v)\vee(u\wedge n(v))=u\wedge v\leq v.
\]

Now set
\[
 t=x\wedge(y\vee z),\qquad
 d=(x\wedge y)\vee(x\wedge z).
\]
Both $x\wedge y$ and $x\wedge z$ lie below $t$, so $d\leq t$.
To prove the other inequality, consider $r=t\wedge n(d)$.
We will show $r=\zero$, which implies $t\leq d$ by
\eqref{eq:orthogonal}.

We have $r\leq x$ and $r\leq n(d)$. Since $x\wedge y\leq d$,
order reversal gives $r\leq n(x\wedge y)$. Consequently,
\[
 r\wedge y\leq x\wedge y,
 \qquad
 r\wedge y\leq r\leq n(x\wedge y).
\]
Thus $r\wedge y$ is below the meet of an element and its complement,
so $r\wedge y=\zero$. By \eqref{eq:orthogonal} and double negation,
$r\leq n(y)$. The same argument with $z$ gives $r\leq n(z)$.
Using \eqref{eq:demorgan},
\[
 r\leq n(y)\wedge n(z)=n(y\vee z).
\]
But also $r\leq t\leq y\vee z$. Therefore $r$ lies below an element
and its complement, so $r=\zero$, as desired.
We have proved
\begin{equation}
 x\wedge(y\vee z)=(x\wedge y)\vee(x\wedge z).\label{eq:distrib}
\end{equation}
Applying $n$ to \eqref{eq:distrib} with $x,y,z$ replaced by
$n(x),n(y),n(z)$, and using De Morgan's laws, gives the dual
distributive law.

It follows that $(A,\wedge,\vee,n,\zero,\one)$ is a Boolean algebra.
Moreover,
\[
 n(x\wedge y)=n(n(x*y))=x*y,
\]
so its NAND operation is exactly the operation with which we started.

\section{Converse and equational completeness}\label{sec:complete}
\subsection{Checking the axiom in a Boolean algebra}
In any Boolean algebra, let $a*b=\neg(a\wedge b)$.
The ordinary Boolean laws give
\[
 (a*b)*c=(a\wedge b)\vee\neg c,
 \qquad
 a*((a*c)*a)=\neg a\vee\neg c.
\]
For the second equality, one can first calculate
$(a*c)*a=(a\wedge c)\vee\neg a$, and then take NAND with $a$.
It follows that
\[
 \begin{aligned}
 &((a*b)*c)*(a*((a*c)*a))\\
 &\qquad=\neg\bigl(((a\wedge b)\vee\neg c)
                    \wedge(\neg a\vee\neg c)\bigr)\\
 &\qquad=\neg\bigl(\neg c\vee(a\wedge b\wedge\neg a)\bigr)\\
 &\qquad=c.
 \end{aligned}
\]
Thus \eqref{eq:W} holds, completing the proof of
Theorem~\ref{thm:main}.\hfill$\square$

\subsection{Why every Boolean NAND identity follows}
The model characterization also gives the claimed completeness. We spell
out the short quotient argument so that this last step does not rely on
an unstated completeness theorem.

Let $T$ be the set of terms formed from variables $x_0,x_1,\ldots$ and a
single binary operation symbol $*$. Write $s\equiv t$ if the equation
$s=t$ can be derived using substitution instances of \eqref{eq:W},
reflexivity, symmetry, transitivity, and replacement of equal subterms.
This is an equivalence relation and a congruence:
\[
 s\equiv s',\quad t\equiv t'
 \quad\Longrightarrow\quad
 s*t\equiv s'*t'.
\]
Thus the quotient $F=T/{\equiv}$ has a well-defined operation
$[s]*[t]=[s*t]$. It is nonempty, since it contains $[x_0]$.
All substitution instances of \eqref{eq:W} are among the generating
relations, so $F$ satisfies \eqref{eq:W}. By
Theorem~\ref{thm:main}, $F$ is the NAND reduct of a Boolean algebra.

Suppose now that $s=t$ is valid under every interpretation in every
Boolean algebra. Interpret each variable $x_i$ as $[x_i]\in F$.
Induction on terms shows that a term $r$ then evaluates to $[r]$.
Validity therefore gives $[s]=[t]$, which, by definition of the quotient,
means $s\equiv t$. The identity is derivable from \eqref{eq:W}.
Conversely, every derivable identity is valid in every Boolean algebra,
because the axiom and all equality rules are valid there. We conclude:
\begin{theorem}[Equational completeness]
An identity in the NAND language holds in every Boolean algebra if and
only if it is derivable from \eqref{eq:W} by equational reasoning.
\end{theorem}

\end{document}
